%%%%%%%%%%%%%%%%%%%%%%%%%%%%%%%%%%%%%%%%%%%%%%%%%%%%%%%%%%%%%%%%%%%%%%
% Lớp 10 | Thống kê | Số gần đúng. Sai số
% Dạng 0D6?1-2: Tính và xác định độ chính xác của kết quả
% Số câu: 7  (sinh tự động bởi tach_cau.py; từ nay thầy sửa trực tiếp file này)
%%%%%%%%%%%%%%%%%%%%%%%%%%%%%%%%%%%%%%%%%%%%%%%%%%%%%%%%%%%%%%%%%%%%%%

%%%%%%%%%%%%%%%% TRẮC NGHIỆM 4 PHƯƠNG ÁN %%%%%%%%%%%%%%%%

% nguồn: MaHoaNH_TracNghiem_Lop10
\begin{ex}%[0D6N1-2]%[ma:4b68fc39]
	Độ cao của một ngọn  núi được ghi lại như sau $\overline{h}=1372{,}5\,\text{m}\pm 0{,}2\,\text{m}$. Độ chính xác $ d$ của phép đo trên là
	\choice 
		{$ d=0{,}1\,\text{m}$} 
		{$ d=1\,\text{m}$} 
		 {\True $ d=0{,}2\,\text{m}$}
		 {$ d=2\,\text{m}$}
	\loigiai{ 
	Độ chính xác $  d=0{,}2\,\text{m}$.}
\end{ex}

% nguồn: MaHoaNH_TracNghiem_Lop10
\begin{ex}%[0D6H1-2]%[ma:36333eab]
	Chiều cao của một ngọn đồi là $\overline{h}=347{,}13$$\pm 0,2$  m. Độ chính xác $ d$ của phép đo trên là
	\choice 
		{$ d=347{,}13$ m} 
		{$ 347{,}33$ m} 
		{\True $ d=0{,}2$  m}
		 {$ d=346{,}93$ m}
	 \loigiai{
	Ta có $ a$ là số gần đúng của $\overline{a}$ với độ chính  xác $ d$ qui ước viết gọn là $\overline{a}=a\pm d$. Vậy độ chính xác của phép đo là $ d=0{,}2$ m.}
\end{ex}

% nguồn: Ok_TracNghiemLop11_Lan1

%%%%%%%%%%%%%%%% ĐÚNG/SAI %%%%%%%%%%%%%%%%

% nguồn: MaHoaNH_DungSai_Lop10
\begin{ex}%[0D6H1-2]%[ma:706a4396]
	Một mảnh đất hình chữ nhật có chiều dài $a= 25\pm 0{,}01$ m, chiều rộng $b= 15\pm  0{,}01$ m.
	Các khẳng định  sau đây đúng hay sai?
	\choiceTF 
		{\True Độ chính  xác phép đo chiều dài là $0{,}01$ m}
		{Sai số tương đối  của phép đo chiều rộng không vượt quá $0{,}05\%$}
		{\True Chu vi của mảnh không vượt quá $80\pm 0{,}04$ m} 
		{\True Diện tích của  mảnh đất không nhỏ hơn $374{,}6001~\mathrm{m^2}$}
	\loigiai{ 
		\begin{itemchoice}
			\itemch Độ chính xác phép đo chiều  dài là $0{,}01$ m.
			\itemch Sai số tương đối của phép đo chiều rộng không vượt quá $\dfrac{0,01}{15} \approx  0,0667\%$.
			\itemch Chu vi của mảnh  không vượt quá $2(25{,}01+15{,}01) = 80\pm 0{,}04$ m.
			 \itemch Diện tích của mảnh đất không nhỏ hơn $24{,}99 \cdot 14{,}99 =374{,}6001 ~\mathrm{m^2}$.
		\end{itemchoice}
	 }
\end{ex}

% nguồn: MaHoaNH_DungSai_Lop10
\begin{ex}%[0D6H1-2]%[ma:de76562e]
	Kết quả một phép đo  được viết là $\ell= 180{,}25 \pm 0{,}005$. \\
	Các  khẳng định sau đây đúng hay sai?
	\choiceTF 
		{Độ chính xác  của phép đo là $\pm 0{,}005$}
		{\True Kết quả của phép đo thuộc đoạn $[180{,}245;  180{,}255]$}
		{\True Sai  số tuyệt đối của phép đo không vượt quá $0{,}005$}
		{Sai số tương đối của phép đo  không vượt quá $0{,}002\%$}
	 \loigiai{
		\begin{itemchoice}
			\itemch Độ chính xác  của phép đo là $0{,}005$.
			\itemch Ta có $180{,}25 - 0{,}005 \leq \ell  \leq 180{,}25 + 0{,}005$. \\
			Kết quả của phép đo thuộc đoạn  $[180{,}245; 180{,}255]$.
			\itemch Ta có $\Delta_a \leq d$ nên sai số tuyệt đối của phép đo không vượt  quá $0{,}005$.
			\itemch Ta có $\delta_\ell = \dfrac{\Delta_\ell}{|\ell|} \leq \dfrac{0{,}005}{180{,}25} \approx  0{,}0028\%$
		\end{itemchoice}
	} 
\end{ex}

%%%%%%%%%%%%%%%% TRẢ LỜI NGẮN %%%%%%%%%%%%%%%%

% nguồn: MaHoaNH_TraLoiNgan_Lop10
\begin{bt}%[0D6H1-2]%[ma:e741074f]
	Bạn Nam thực hiện việc  đo đường kính của một ống đồng bằng thước kẹp. Kết quả của Nam là $15{,}3$ mm. Hỏi số số gần đúng của phép đo là bao nhiêu?
	\shortans{15,3} 
	 \loigiai{
	Số $15{,}3$ là số gần đúng của  phép đo trong thực tế}
\end{bt}

% nguồn: MaHoaNH_TraLoiNgan_Lop10
\begin{bt}%[0D6H1-2]%[ma:015c1fe0]
	Trên bao bì của một sản phẩm có ghi \lq\lq khối lượng tịnh $200\pm 2$ g\rq\rq. Hãy đánh giá sai số tương đối của  số gần đúng này.
	\shortans{1} 
	 \loigiai{
		Khối lượng  đúng của bao sản phẩm $\bar{a}$ (tính theo gam) thuộc đoạn $[198; 202]$. \\
	Sai số tương  đối là $\delta_a\leq\dfrac{2}{200}=1\%$.}
\end{bt}

% nguồn: TuLuan_Lop10_daduyet (nhập 14/11/2026)
\begin{bt}[H]%[Nguyễn Văn Cường (Cường NV) - Dự án bài giảng K10 - K11]%[0D6H1-2]%[ma:3397e7ae]
Diện tích hình tròn bán kính $r = 12$ cm là $S = \pi \cdot r^2 \approx 3{,}145 \cdot 12^2 \approx$ $452{,}88$ cm$^2$. Biết rằng $3{,}14<\pi<3{,}15$. Hãy ước lượng độ chính xác của $S$.
\loigiai{
Gọi $\bar{S}$ là diện tích đúng của hình tròn, ta có
\[3{,}14 \cdot 12^2 < \bar{S} <3{,}15 \cdot 12^2.\]
Hay $452{,}16 < \bar{S} < 453{,}6$.\\
Suy ra $452{,}16 - 452{,}88 < \bar{S}- S <453{,}6-452{,}88\Rightarrow |\bar{S} - S| < 0{,}72$.\\
Vậy kết quả có độ chính xác là $0{,}72$ hay diện tích của hình tròn là $452{,}88 \pm 0{,}72$ cm$^2$.}
\end{bt}

% nguồn: TuLuan_Lop10_Dot2 (nhập 14/11/2026)
\begin{bt}%[0D6H1-2]%[ma:5b1f2ad7]
	Một người thợ cần biết chiều cao của một ngôi nhà. Anh ta thực hiện các phép đo trong ba lần và được kết quả như sau: $ h_1=10{,}23\pm0{,}43 $ (m), $ h_2=10{,}58\pm0{,}2 $ (m), $ h_3=9{,}92\pm0{,}63 $ (m). Hỏi trong ba số liệu đó, người thợ nên chọn số nào làm chiều cao ngôi nhà.
	\loigiai{
		Phép đo lần 1 có sai số tương đối $ \delta_1\leq\dfrac{0{,}43}{10{,}23} \approx0{,}042=4{,}2\% $.\\
		Phép đo lần 2 có sai số tương đối $ \delta_2\leq\dfrac{0{,}2}{10{,}58} \approx0{,}0189=1{,}89\% $.\\
		Phép đo lần 3 có sai số tương đối $ \delta_3\leq\dfrac{0{,}63}{9{,}92} \approx0{,}0635=6{,}35\% $.\\
		Như vậy người thợ nên chọn $ h_2=10{,}58\pm0{,}2 $ (m) làm chiều cao ngôi nhàn.
	} 
\end{bt}

% nguồn: ThucTe_Lop10_TLN_DS (nhập 14/11/2026)
\begin{ex}%[0D6H1-2][H]%[Bài giảng 10 New - 4in1, Nguyễn Vân Trường]%[ma:f613db51]
	Độ dài của một đoạn thẳng người ta đo được là $ 2 $ m $ \pm$ $0{,}05 $ m (làm tròn đến phần trăm).
	\choiceTF[t]
	{Độ dài gần đúng của đoạn thẳng là $ 2{,}05 $ m}
	{\True Độ chính xác của phép đo là $ d=0{,}05 $}
	{Sai số  tuyệt đối $ \Delta_a\leq d=0{,}05 $}
	{\True Sai số tương đối tối đa cho phép trong phép đo trên là $ 0{,}03\% $}
	\loigiai{
			\begin{itemchoice}
		\itemch Ta có độ dài gần đúng của đoạn thẳng là $ a=2 $.
		\itemch Độ chính xác là $ d=0{,}05 $.\\
		\itemch Vì sai số  tuyệt đối $ \Delta_a\leq d=0{,}05 $ nên sai số tương đối là $ \delta_a\leq \dfrac{d}{\lvert a\rvert}=\dfrac{0{,}05}{2}\approx 0{,}03\% $.\\
		\itemch Vậy sai số tương đối tối đa cho phép trong phép đo trên là $ 0{,}03\% $.
	\end{itemchoice}
	}
\end{ex}
